\chapter{Simulados Completos}

\section{Como usar}

O simulado reproduz a arquitetura objetiva do edital: \textbf{100 questões de múltipla escolha}, cada uma com cinco alternativas e uma resposta correta, divididas em \textbf{40 questões de conhecimentos gerais} e \textbf{60 de conhecimentos específicos}. Para aproximar a pontuação real, considere 1 ponto por questão geral e 2 pontos por questão específica, totalizando 160 pontos.

\begin{atencao}
Não consulte o gabarito durante a prova. Faça o bloco objetivo em até 5 horas, marque dúvidas e somente depois corrija. Classifique cada erro como: falta de teoria, interpretação, cálculo, confusão entre conceitos ou distração.
\end{atencao}

\section{Simulado 1 — 100 questões}

\subsection*{Conhecimentos Gerais — questões 1 a 40}

\begin{questao}{S1.01 — Língua Portuguesa}
Assinale a reescrita que preserva o sentido de: ``Embora o sistema tenha sido atualizado, persistiram falhas que impediram a conclusão da auditoria.''
\begin{enumerate}[label=\Alph*)]
\item Como o sistema foi atualizado, desapareceram as falhas.
\item Apesar de o sistema ter sido atualizado, permaneceram falhas que inviabilizaram a conclusão da auditoria.
\item O sistema foi atualizado porque as falhas impediram a auditoria.
\item As falhas persistiram, portanto o sistema não foi atualizado.
\item A auditoria foi concluída, ainda que houvesse falhas.
\end{enumerate}
\end{questao}

\begin{questao}{S1.02 — Língua Portuguesa}
Em ``Os relatórios que apresentavam inconsistências foram devolvidos'', a ausência de vírgulas indica que:
\begin{enumerate}[label=\Alph*)]
\item todos os relatórios apresentavam inconsistências.
\item apenas o subconjunto de relatórios que apresentava inconsistências foi devolvido.
\item nenhum relatório foi devolvido.
\item a oração é necessariamente explicativa.
\item a pontuação é indiferente ao sentido.
\end{enumerate}
\end{questao}

\begin{questao}{S1.03 — Língua Portuguesa}
Assinale a opção em que a concordância está de acordo com a norma-padrão.
\begin{enumerate}[label=\Alph*)]
\item Houveram falhas relevantes no processo.
\item Fazem três meses que a auditoria começou.
\item Devem existir evidências suficientes para a conclusão.
\item Existe, nos autos, documentos contraditórios.
\item Tratam-se de controles essenciais.
\end{enumerate}
\end{questao}

\begin{questao}{S1.04 — Língua Portuguesa}
A substituição de ``porquanto'' por qual conectivo preserva, em regra, valor causal/explicativo?
\begin{enumerate}[label=\Alph*)]
\item contudo.
\item porque.
\item portanto.
\item embora.
\item à medida que.
\end{enumerate}
\end{questao}

\begin{questao}{S1.05 — Língua Portuguesa}
Em ``A equipe informou ao gestor que seu relatório continha inconsistências'', há potencial ambiguidade porque o pronome ``seu'' pode retomar:
\begin{enumerate}[label=\Alph*)]
\item apenas equipe.
\item apenas gestor.
\item equipe ou gestor, conforme o contexto.
\item inconsistências.
\item verbo informar.
\end{enumerate}
\end{questao}

\begin{questao}{S1.06 — Língua Portuguesa}
Assinale a opção em que o emprego da crase está correto.
\begin{enumerate}[label=\Alph*)]
\item A auditoria foi realizada à partir de maio.
\item O relatório fez referência à norma vigente.
\item A equipe entregou o documento à todos.
\item O auditor ficou frente à frente com o gestor.
\item O sistema funciona de segunda à sexta.
\end{enumerate}
\end{questao}

\begin{questao}{S1.07 — Competências Digitais}
Um servidor recebe e-mail com link para página visualmente idêntica ao portal institucional, mas o domínio possui pequena alteração ortográfica. A ação inicial mais segura é:
\begin{enumerate}[label=\Alph*)]
\item inserir a senha e verificar depois.
\item ignorar o domínio porque a aparência é confiável.
\item não acessar pelo link, validar o endereço por canal confiável e reportar possível phishing.
\item encaminhar a mensagem a todos os colegas.
\item desativar o antivírus para evitar bloqueio falso positivo.
\end{enumerate}
\end{questao}

\begin{questao}{S1.08 — Competências Digitais}
Segundo a LGPD, o princípio da necessidade orienta que o tratamento de dados pessoais:
\begin{enumerate}[label=\Alph*)]
\item colete todo dado disponível para uso futuro indeterminado.
\item se limite ao mínimo necessário para a finalidade legítima informada.
\item seja sempre baseado em consentimento.
\item dispense medidas de segurança.
\item exclua qualquer tratamento pelo poder público.
\end{enumerate}
\end{questao}

\begin{questao}{S1.09 — Competências Digitais}
A LAI parte da lógica de que:
\begin{enumerate}[label=\Alph*)]
\item sigilo é regra e publicidade é exceção.
\item publicidade é regra e sigilo é exceção, observadas hipóteses legais.
\item toda informação pessoal é pública.
\item pedidos de acesso precisam demonstrar interesse jurídico específico.
\item somente o Judiciário pode fornecer informação pública.
\end{enumerate}
\end{questao}

\begin{questao}{S1.10 — Competências Digitais}
No Marco Civil da Internet, registros e dados devem ser tratados segundo regras legais de privacidade e proteção, o que significa que:
\begin{enumerate}[label=\Alph*)]
\item qualquer servidor pode consultar dados sem finalidade definida.
\item segurança, finalidade e acesso autorizado continuam relevantes no setor público.
\item a internet não se submete à proteção de dados.
\item logs devem ser publicados integralmente.
\item dados de terceiros podem ser compartilhados livremente entre órgãos.
\end{enumerate}
\end{questao}

\begin{questao}{S1.11 — Raciocínio Lógico}
A negação lógica de ``Todos os relatórios foram aprovados'' é:
\begin{enumerate}[label=\Alph*)]
\item Nenhum relatório foi aprovado.
\item Pelo menos um relatório não foi aprovado.
\item Todos os relatórios foram reprovados.
\item Alguns relatórios foram aprovados.
\item Exatamente um relatório foi reprovado.
\end{enumerate}
\end{questao}

\begin{questao}{S1.12 — Raciocínio Lógico}
A proposição ``Se o controle é eficaz, então o risco é reduzido'' é logicamente equivalente a:
\begin{enumerate}[label=\Alph*)]
\item Se o risco não é reduzido, então o controle não é eficaz.
\item Se o risco é reduzido, então o controle é eficaz.
\item O controle é eficaz e o risco não é reduzido.
\item O controle não é eficaz e o risco é reduzido.
\item O risco é reduzido se e somente se o controle é eficaz.
\end{enumerate}
\end{questao}

\begin{questao}{S1.13 — Raciocínio Lógico}
Um orçamento de R\$ 800 mil sofreu acréscimo de 15\%. O novo valor é:
\begin{enumerate}[label=\Alph*)]
\item R\$ 880 mil.
\item R\$ 900 mil.
\item R\$ 920 mil.
\item R\$ 940 mil.
\item R\$ 960 mil.
\end{enumerate}
\end{questao}

\begin{questao}{S1.14 — Raciocínio Lógico}
Em um conjunto de 200 processos, 120 possuem falha A, 80 possuem falha B e 40 possuem ambas. Quantos possuem ao menos uma das duas falhas?
\begin{enumerate}[label=\Alph*)]
\item 120.
\item 140.
\item 160.
\item 180.
\item 200.
\end{enumerate}
\end{questao}

\begin{questao}{S1.15 — Raciocínio Lógico}
Se $x+2y=14$ e $x-y=2$, então $x$ é igual a:
\begin{enumerate}[label=\Alph*)]
\item 4.
\item 5.
\item 6.
\item 7.
\item 8.
\end{enumerate}
\end{questao}

\begin{questao}{S1.16 — Direito Administrativo}
A desconcentração administrativa ocorre quando:
\begin{enumerate}[label=\Alph*)]
\item há criação de nova pessoa jurídica.
\item competências são distribuídas internamente dentro da mesma pessoa jurídica.
\item serviço público é necessariamente privatizado.
\item competência é transferida para outro ente federativo por emenda constitucional.
\item autarquia é extinta.
\end{enumerate}
\end{questao}

\begin{questao}{S1.17 — Direito Administrativo}
O atributo do ato administrativo que permite, em certas hipóteses legais, execução direta pela Administração sem ordem judicial é:
\begin{enumerate}[label=\Alph*)]
\item tipicidade.
\item presunção de legitimidade.
\item autoexecutoriedade.
\item competência.
\item finalidade.
\end{enumerate}
\end{questao}

\begin{questao}{S1.18 — Direito Administrativo}
A anulação de ato administrativo decorre, em regra, de:
\begin{enumerate}[label=\Alph*)]
\item conveniência e oportunidade.
\item ilegalidade.
\item mudança política sem vício.
\item simples desejo do particular.
\item decurso de prazo contratual.
\end{enumerate}
\end{questao}

\begin{questao}{S1.19 — Direito Administrativo}
Após as alterações da Lei de Improbidade Administrativa, é correto afirmar que:
\begin{enumerate}[label=\Alph*)]
\item qualquer culpa caracteriza improbidade.
\item improbidade exige enquadramento legal e dolo nos termos da disciplina vigente.
\item toda ilegalidade administrativa é improbidade.
\item sanção de improbidade dispensa devido processo.
\item dano ao erário nunca integra improbidade.
\end{enumerate}
\end{questao}

\begin{questao}{S1.20 — Direito Administrativo}
A responsabilidade civil do Estado por dano causado por agente público, nessa qualidade, é em regra:
\begin{enumerate}[label=\Alph*)]
\item subjetiva para a vítima, exigindo prova de dolo do agente.
\item objetiva quanto ao Estado, sem eliminar análise de nexo e excludentes aplicáveis.
\item inexistente.
\item exclusivamente penal.
\item sempre solidária e automática com o agente.
\end{enumerate}
\end{questao}

\begin{questao}{S1.21 — Direito Constitucional}
No federalismo brasileiro, os municípios:
\begin{enumerate}[label=\Alph*)]
\item não integram a organização político-administrativa.
\item integram a Federação com autonomia nos limites constitucionais.
\item são subordinados hierarquicamente aos estados.
\item podem se separar unilateralmente da Federação.
\item exercem competências exclusivas da União.
\end{enumerate}
\end{questao}

\begin{questao}{S1.22 — Direito Constitucional}
Os direitos fundamentais previstos na Constituição:
\begin{enumerate}[label=\Alph*)]
\item restringem-se aos direitos individuais do art. 5º.
\item incluem também direitos sociais, políticos e outros direitos constitucionalmente protegidos.
\item são sempre absolutos.
\item não vinculam o poder público.
\item não admitem proteção judicial.
\end{enumerate}
\end{questao}

\begin{questao}{S1.23 — Direito Constitucional}
A fiscalização contábil, financeira, orçamentária, operacional e patrimonial da União é exercida pelo Congresso Nacional mediante controle externo, com auxílio:
\begin{enumerate}[label=\Alph*)]
\item do STF.
\item do TCU.
\item do CNJ.
\item do Ministério da Fazenda apenas.
\item da Defensoria Pública.
\end{enumerate}
\end{questao}

\begin{questao}{S1.24 — Direito Constitucional}
Comissão parlamentar de inquérito possui poderes de investigação próprios das autoridades judiciais, mas:
\begin{enumerate}[label=\Alph*)]
\item pode condenar criminalmente.
\item não pode praticar atos submetidos à reserva jurisdicional.
\item substitui definitivamente o Poder Judiciário.
\item pode decretar qualquer prisão como pena.
\item pode revogar sentença judicial transitada em julgado.
\end{enumerate}
\end{questao}

\begin{questao}{S1.25 — Direito Constitucional}
O Ministério Público é classificado constitucionalmente como:
\begin{enumerate}[label=\Alph*)]
\item órgão do Poder Judiciário.
\item função essencial à Justiça.
\item órgão auxiliar do Executivo.
\item empresa pública.
\item autarquia especial.
\end{enumerate}
\end{questao}

\begin{questao}{S1.26 — Controle Externo}
Controle externo exercido por Tribunal de Contas distingue-se do controle interno porque:
\begin{enumerate}[label=\Alph*)]
\item apenas o controle interno examina legalidade.
\item possuem posição institucional e funções distintas, embora sejam complementares.
\item controle externo pertence à própria unidade auditada.
\item controle interno é exercido exclusivamente pelo Judiciário.
\item não existe cooperação entre os sistemas.
\end{enumerate}
\end{questao}

\begin{questao}{S1.27 — Controle Externo}
A auditoria operacional concentra-se especialmente em:
\begin{enumerate}[label=\Alph*)]
\item apenas forma contábil.
\item economicidade, eficiência, eficácia e efetividade de políticas e organizações.
\item julgamento penal.
\item controle de constitucionalidade abstrato.
\item registro civil.
\end{enumerate}
\end{questao}

\begin{questao}{S1.28 — Controle Externo}
No sistema brasileiro de jurisdição una:
\begin{enumerate}[label=\Alph*)]
\item decisões administrativas jamais podem ser apreciadas judicialmente.
\item lesão ou ameaça a direito pode ser submetida ao Poder Judiciário.
\item existe contencioso administrativo com coisa julgada judicial obrigatória.
\item Tribunal de Contas integra o Judiciário.
\item Administração não possui processos próprios.
\end{enumerate}
\end{questao}

\begin{questao}{S1.29 — Controle Externo}
O parecer prévio sobre contas de governo:
\begin{enumerate}[label=\Alph*)]
\item é idêntico a sentença penal.
\item integra o modelo de apreciação das contas do chefe do Executivo e subsidia julgamento político competente.
\item elimina participação do Legislativo.
\item é sempre produzido por unidade técnica sem deliberação do Tribunal.
\item equivale necessariamente à imputação de débito.
\end{enumerate}
\end{questao}

\begin{questao}{S1.30 — Legislação TCE-MA}
A IN TCE-MA nº 50/2017 relaciona-se principalmente a:
\begin{enumerate}[label=\Alph*)]
\item tomada de contas especial.
\item emendas parlamentares impositivas.
\item contratação de nuvem.
\item acessibilidade digital.
\item organização de partidos políticos.
\end{enumerate}
\end{questao}

\begin{questao}{S1.31 — Legislação TCE-MA}
A IN TCE-MA nº 82/2025 possui foco em:
\begin{enumerate}[label=\Alph*)]
\item execução de emendas parlamentares, transparência e rastreabilidade.
\item processo eleitoral interno do Tribunal.
\item gestão de folha de pagamento federal.
\item tributação estadual.
\item licenciamento ambiental de mineração.
\end{enumerate}
\end{questao}

\begin{questao}{S1.32 — Legislação TCE-MA}
Regimento Interno e Lei Orgânica do TCE-MA devem ser interpretados de modo que:
\begin{enumerate}[label=\Alph*)]
\item o Regimento possa contrariar livremente a lei.
\item a lei prevaleça sobre regra regimental incompatível.
\item ambos tenham exatamente a mesma hierarquia formal.
\item norma técnica interna revogue lei automaticamente.
\item costume administrativo prevaleça sempre.
\end{enumerate}
\end{questao}

\begin{questao}{S1.33 — Legislação TCE-MA}
Em matéria de emendas parlamentares, rastreabilidade significa:
\begin{enumerate}[label=\Alph*)]
\item divulgar apenas o valor total anual.
\item permitir reconstruir proponente, finalidade, execução financeira, documentos e resultado.
\item ocultar o beneficiário.
\item excluir processo administrativo da cadeia de informação.
\item guardar dados somente em planilha privada.
\end{enumerate}
\end{questao}

\begin{questao}{S1.34 — História/Geografia MA}
A Batalha de Guaxenduba está associada:
\begin{enumerate}[label=\Alph*)]
\item à reação contra a presença francesa.
\item à Balaiada.
\item à Revolução de 1930.
\item à invasão holandesa de 1641.
\item à Greve de 1951.
\end{enumerate}
\end{questao}

\begin{questao}{S1.35 — História/Geografia MA}
A Balaiada ocorreu:
\begin{enumerate}[label=\Alph*)]
\item no Período Regencial, com expressiva participação popular.
\item durante a França Equinocial.
\item apenas em São Luís em 1951.
\item no Estado Novo.
\item após 1964.
\end{enumerate}
\end{questao}

\begin{questao}{S1.36 — História/Geografia MA}
Qual conjunto reúne rios/bacias destacados no edital como limítrofes?
\begin{enumerate}[label=\Alph*)]
\item Mearim, Itapecuru e Munim.
\item Parnaíba, Gurupi e Tocantins-Araguaia.
\item Amazonas, São Francisco e Paraná.
\item Negro, Solimões e Madeira.
\item Jaguaribe, Capibaribe e Iguaçu.
\end{enumerate}
\end{questao}

\begin{questao}{S1.37 — História/Geografia MA}
A importância logística do Maranhão relaciona-se, entre outros fatores:
\begin{enumerate}[label=\Alph*)]
\item ao Porto do Itaqui e a corredores ferroviários de exportação.
\item à ausência de acesso marítimo.
\item à inexistência de agronegócio.
\item à proibição de transporte ferroviário.
\item ao isolamento de São Luís de cadeias nacionais.
\end{enumerate}
\end{questao}

\begin{questao}{S1.38 — Direitos Humanos}
Segundo a LBI, a deficiência:
\begin{enumerate}[label=\Alph*)]
\item gera incapacidade civil automática.
\item é analisada na interação entre impedimentos e barreiras à participação.
\item é definida apenas por renda.
\item impede exercício de direitos familiares.
\item depende sempre de decisão judicial.
\end{enumerate}
\end{questao}

\begin{questao}{S1.39 — Direitos Humanos}
O ODS 16 está diretamente associado a:
\begin{enumerate}[label=\Alph*)]
\item instituições eficazes, responsáveis e inclusivas, paz e justiça.
\item apenas produção agrícola.
\item exclusivamente energia.
\item somente oceanos.
\item apenas turismo.
\end{enumerate}
\end{questao}

\begin{questao}{S1.40 — Direitos Humanos}
Ações afirmativas são compatíveis com igualdade material quando:
\begin{enumerate}[label=\Alph*)]
\item buscam corrigir desigualdades e promover igualdade de oportunidades.
\item criam distinção arbitrária sem finalidade legítima.
\item eliminam direitos fundamentais.
\item substituem toda política universal.
\item dispensam qualquer base normativa.
\end{enumerate}
\end{questao}

\subsection*{Conhecimentos Específicos — questões 41 a 100}

\begin{questao}{S1.41 — Infraestrutura de TI}
Uma rede \texttt{10.0.0.0/26} possui, no bloco, quantos endereços IPv4?
\begin{enumerate}[label=\Alph*)]\item 30.\item 32.\item 62.\item 64.\item 128.\end{enumerate}
\end{questao}

\begin{questao}{S1.42 — Infraestrutura de TI}
Para autenticação centralizada de acesso a portas de rede cabeada e Wi-Fi corporativo, a combinação mais apropriada é:
\begin{enumerate}[label=\Alph*)]\item 802.1X, EAP e RADIUS.\item SNMP e NAT.\item RAID e DNS.\item MPLS e SMTP.\item NTP e FTP.\end{enumerate}
\end{questao}

\begin{questao}{S1.43 — Infraestrutura de TI}
Um sistema mantém sessões apenas na memória local de cada servidor. Após failover, usuários perdem autenticação. A melhoria arquitetural é:
\begin{enumerate}[label=\Alph*)]\item externalizar/replicar estado ou adotar sessão stateless.\item trocar TCP por UDP.\item usar RAID 0.\item aumentar TTL do DNS.\item remover health checks.\end{enumerate}
\end{questao}

\begin{questao}{S1.44 — Infraestrutura de TI}
RAID 6 e snapshots no mesmo storage não substituem backup externo porque:
\begin{enumerate}[label=\Alph*)]\item não protegem adequadamente contra todos os eventos do mesmo domínio de falha, corrupção e ransomware.\item RAID 6 não possui paridade.\item snapshots sempre são offline.\item backup serve apenas para redes.\item RAID 6 suporta somente um disco.\end{enumerate}
\end{questao}

\begin{questao}{S1.45 — Infraestrutura de TI}
As quatro condições clássicas necessárias ao deadlock incluem:
\begin{enumerate}[label=\Alph*)]\item exclusão mútua, posse e espera, ausência de preempção e espera circular.\item apenas alta CPU.\item fragmentação, NAT, DNS e ARP.\item cache, TLB, swap e RAID.\item autenticação, autorização, auditoria e accounting.\end{enumerate}
\end{questao}

\begin{questao}{S1.46 — Infraestrutura de TI}
No Active Directory, é correto afirmar que:
\begin{enumerate}[label=\Alph*)]\item LDAP e AD são sinônimos perfeitos.\item AD utiliza, entre outros, LDAP, Kerberos e DNS.\item Kerberos não usa tickets.\item GPO não permite configuração centralizada.\item domínio e floresta são sempre idênticos.\end{enumerate}
\end{questao}

\begin{questao}{S1.47 — Engenharia de Dados}
Uma tabela em que atributo não chave depende de outro atributo não chave viola, tipicamente:
\begin{enumerate}[label=\Alph*)]\item 3FN.\item atomicidade.\item isolamento de transação.\item integridade física do disco.\item CAP.\end{enumerate}
\end{questao}

\begin{questao}{S1.48 — Engenharia de Dados}
Em banco com MVCC, uma finalidade central é:
\begin{enumerate}[label=\Alph*)]\item permitir versões para reduzir bloqueios entre leitura e escrita conforme o isolamento.\item eliminar transações.\item impedir índices.\item substituir logs de recuperação.\item eliminar consistência.\end{enumerate}
\end{questao}

\begin{questao}{S1.49 — Engenharia de Dados}
Em modelo dimensional, tabela fato contém principalmente:
\begin{enumerate}[label=\Alph*)]\item medidas de eventos e chaves para dimensões.\item apenas descrições textuais.\item somente usuários do SGBD.\item código-fonte de ETL.\item políticas de firewall.\end{enumerate}
\end{questao}

\begin{questao}{S1.50 — Engenharia de Dados}
CDC é técnica utilizada para:
\begin{enumerate}[label=\Alph*)]\item capturar alterações de dados e propagá-las a outros sistemas/pipelines.\item criptografar disco inteiro.\item resolver deadlock em rede.\item executar container.\item classificar texto.\end{enumerate}
\end{questao}

\begin{questao}{S1.51 — Engenharia de Dados}
Lineage é importante porque permite:
\begin{enumerate}[label=\Alph*)]\item rastrear origem, transformações e destino dos dados.\item aumentar clock da CPU.\item eliminar metadados.\item substituir backup.\item gerar senha aleatória.\end{enumerate}
\end{questao}

\begin{questao}{S1.52 — Engenharia de Dados}
No teorema CAP, diante de partição de rede, o sistema distribuído precisa administrar trade-off entre:
\begin{enumerate}[label=\Alph*)]\item consistência e disponibilidade.\item CPU e memória.\item integridade e confidencialidade.\item backup e RAID.\item SQL e HTML.\end{enumerate}
\end{questao}

\begin{questao}{S1.53 — Engenharia de Software}
Um requisito ``o sistema deve ser rápido'' é fraco porque:
\begin{enumerate}[label=\Alph*)]\item não define critério mensurável de aceitação/desempenho.\item todo requisito deve ser código.\item requisitos não funcionais são proibidos.\item rapidez nunca pode ser testada.\item apenas usuário pode medir latência.\end{enumerate}
\end{questao}

\begin{questao}{S1.54 — Engenharia de Software}
No Scrum, a Sprint Review tem como foco principal:
\begin{enumerate}[label=\Alph*)]\item inspecionar incremento e adaptar Product Backlog com stakeholders.\item avaliar individualmente funcionários.\item substituir Daily Scrum.\item aprovar folha de pagamento.\item configurar firewall.\end{enumerate}
\end{questao}

\begin{questao}{S1.55 — Engenharia de Software}
TDD caracteriza-se pelo ciclo simplificado:
\begin{enumerate}[label=\Alph*)]\item teste falha, código mínimo passa, refatoração.\item código completo, depois documentação, sem teste.\item deploy direto em produção.\item especificação informal sem automação.\item pentest antes de requisito.\end{enumerate}
\end{questao}

\begin{questao}{S1.56 — Engenharia de Software}
Uma pipeline CI/CD com credencial administrativa permanente de produção viola principalmente:
\begin{enumerate}[label=\Alph*)]\item princípio do menor privilégio e segregação.\item normalização de dados.\item CAP.\item encapsulamento TCP.\item semântica REST.\end{enumerate}
\end{questao}

\begin{questao}{S1.57 — Engenharia de Software}
Em Kubernetes, readiness probe serve principalmente para:
\begin{enumerate}[label=\Alph*)]\item indicar se o pod está pronto para receber tráfego.\item criptografar secrets.\item criar imagem Docker.\item controlar DNS externo.\item assinar commits Git.\end{enumerate}
\end{questao}

\begin{questao}{S1.58 — Engenharia de Software}
Um SLO de disponibilidade deve ser acompanhado por:
\begin{enumerate}[label=\Alph*)]\item indicador mensurável e janela de avaliação.\item apenas opinião do usuário.\item código sem telemetria.\item ausência de logs.\item métrica sem unidade.\end{enumerate}
\end{questao}

\begin{questao}{S1.59 — Segurança da Informação}
MFA fatigue busca explorar:
\begin{enumerate}[label=\Alph*)]\item repetidas solicitações de autenticação para induzir aprovação indevida.\item colisão de hash por força bruta exclusivamente.\item falha de RAID.\item fragmentação IP.\item compressão de logs.\end{enumerate}
\end{questao}

\begin{questao}{S1.60 — Segurança da Informação}
PAM é mais diretamente associado a:
\begin{enumerate}[label=\Alph*)]\item gestão e controle de acessos privilegiados.\item roteamento BGP.\item criação de data warehouse.\item gerenciamento de projetos.\item compressão de backup.\end{enumerate}
\end{questao}

\begin{questao}{S1.61 — Segurança da Informação}
Assinatura digital oferece principalmente:
\begin{enumerate}[label=\Alph*)]\item integridade, autenticidade e suporte a não repúdio, conforme contexto criptográfico.\item confidencialidade automática do conteúdo.\item anonimato absoluto.\item disponibilidade do servidor.\item backup da chave privada.\end{enumerate}
\end{questao}

\begin{questao}{S1.62 — Segurança da Informação}
Em resposta a ransomware, uma cópia imutável/offline é importante porque:
\begin{enumerate}[label=\Alph*)]\item reduz chance de comprometimento simultâneo pelo atacante.\item elimina necessidade de restauração.\item substitui MFA.\item impede qualquer perda lógica.\item dispensa teste de backup.\end{enumerate}
\end{questao}

\begin{questao}{S1.63 — Segurança da Informação}
Cadeia de custódia de evidência digital busca registrar:
\begin{enumerate}[label=\Alph*)]\item coleta, posse, transferências e integridade da evidência.\item apenas tamanho do arquivo.\item somente endereço IP.\item apenas nome do auditor.\item exclusivamente data do relatório.\end{enumerate}
\end{questao}

\begin{questao}{S1.64 — Segurança da Informação}
Zero Trust recomenda:
\begin{enumerate}[label=\Alph*)]\item verificação contínua, identidade/contexto e menor privilégio.\item confiança automática por estar na rede interna.\item eliminar autenticação.\item liberar movimento lateral.\item usar conta compartilhada de administrador.\end{enumerate}
\end{questao}

\begin{questao}{S1.65 — Gestão e Governança de TI}
No COBIT, governança é associada ao domínio:
\begin{enumerate}[label=\Alph*)]\item EDM.\item BAI apenas.\item DSS apenas.\item MEA apenas.\item APO exclusivamente.\end{enumerate}
\end{questao}

\begin{questao}{S1.66 — Gestão e Governança de TI}
Segundo a distinção conceitual de governança e gestão, governança tende a:
\begin{enumerate}[label=\Alph*)]\item avaliar necessidades, direcionar e monitorar.\item executar exclusivamente atividades operacionais.\item substituir todos os gestores.\item eliminar prestação de contas.\item restringir-se a suporte técnico.\end{enumerate}
\end{questao}

\begin{questao}{S1.67 — Gestão e Governança de TI}
No ITIL 4, valor é cocriado por meio de:
\begin{enumerate}[label=\Alph*)]\item relações de serviço e sistema de valor, não apenas entrega unilateral do provedor.\item hardware isolado.\item SLA sem usuário.\item catálogo sem serviço.\item incidentes sem processo.\end{enumerate}
\end{questao}

\begin{questao}{S1.68 — Gestão e Governança de TI}
Em gerenciamento de riscos, risco é normalmente analisado por:
\begin{enumerate}[label=\Alph*)]\item probabilidade e impacto, considerando contexto e controles.\item apenas custo histórico.\item quantidade de servidores.\item cor do dashboard.\item número de reuniões.\end{enumerate}
\end{questao}

\begin{questao}{S1.69 — Gestão e Governança de TI}
No BPMN, gateway é usado para:
\begin{enumerate}[label=\Alph*)]\item controlar divergência/convergência de fluxo conforme regras.\item armazenar dados em disco.\item autenticar usuário.\item compilar código.\item criptografar mensagem.\end{enumerate}
\end{questao}

\begin{questao}{S1.70 — Gestão e Governança de TI}
Um KPI que melhora enquanto resultado final do serviço piora pode indicar:
\begin{enumerate}[label=\Alph*)]\item métrica local mal alinhada ao objetivo de negócio/serviço.\item garantia de governança madura.\item ausência de risco.\item necessidade de eliminar métricas.\item prova de eficácia global.\end{enumerate}
\end{questao}

\begin{questao}{S1.71 — Fiscalização de Contratos de TI}
Na IN SGD/ME nº 94/2022, o macroprocesso de contratação de TIC é estruturado em:
\begin{enumerate}[label=\Alph*)]\item Planejamento da Contratação, Seleção do Fornecedor e Gestão do Contrato.\item orçamento, empenho e pagamento apenas.\item desenvolvimento, teste e deploy.\item edital, recurso e sentença.\item ETP, nota fiscal e multa apenas.\end{enumerate}
\end{questao}

\begin{questao}{S1.72 — Fiscalização de Contratos de TI}
Pesquisa de preços baseada apenas em três cotações 40\% acima de contratações públicas comparáveis, ignoradas sem justificativa, apresenta risco de:
\begin{enumerate}[label=\Alph*)]\item estimativa inadequada e sobrepreço.\item deadlock.\item spoofing.\item perda de pacote.\item fragmentação de índice.\end{enumerate}
\end{questao}

\begin{questao}{S1.73 — Fiscalização de Contratos de TI}
Contrato mede disponibilidade exclusivamente por planilha produzida pela contratada. O principal problema é:
\begin{enumerate}[label=\Alph*)]\item baixa independência e verificabilidade da evidência de medição.\item excesso de normalização.\item ausência de IPv6.\item falta de RAID.\item uso de API REST.\end{enumerate}
\end{questao}

\begin{questao}{S1.74 — Fiscalização de Contratos de TI}
Lock-in deve ser tratado no planejamento mediante:
\begin{enumerate}[label=\Alph*)]\item portabilidade, interoperabilidade, documentação e estratégia de saída.\item exclusão de backups.\item conta administrativa compartilhada.\item formato proprietário obrigatório sem exportação.\item renovação automática sem análise.\end{enumerate}
\end{questao}

\begin{questao}{S1.75 — Fiscalização de Contratos de TI}
SLA de 99,9\% em mês de 30 dias, medido 24x7, admite aproximadamente quanto de indisponibilidade teórica?
\begin{enumerate}[label=\Alph*)]\item 4,32 min.\item 43,2 min.\item 72 min.\item 432 min.\item 720 min.\end{enumerate}
\end{questao}

\begin{questao}{S1.76 — Fiscalização de Contratos de TI}
Quantitativo de 10 mil licenças baseado apenas no total de servidores, quando somente 3.800 usuários utilizam solução equivalente, indica risco de:
\begin{enumerate}[label=\Alph*)]\item superdimensionamento e ociosidade.\item sub-rede pequena.\item ataque DDoS.\item alta cardinalidade.\item ausência de checksum.\end{enumerate}
\end{questao}

\begin{questao}{S1.77 — Computação em Nuvem}
No modelo IaaS, o cliente normalmente é responsável por:
\begin{enumerate}[label=\Alph*)]\item sistema operacional e aplicações das VMs, além de configurações sob sua camada.\item segurança física do data center do provedor.\item fabricação dos discos.\item roteamento global da internet.\item toda a infraestrutura física.\end{enumerate}
\end{questao}

\begin{questao}{S1.78 — Computação em Nuvem}
Elasticidade difere de escalabilidade porque elasticidade enfatiza:
\begin{enumerate}[label=\Alph*)]\item ajuste dinâmico de recursos conforme demanda.\item uso exclusivo de hardware local.\item ausência de automação.\item backup offline.\item criptografia assimétrica.\end{enumerate}
\end{questao}

\begin{questao}{S1.79 — Computação em Nuvem}
Shared responsibility model significa que:
\begin{enumerate}[label=\Alph*)]\item responsabilidades de segurança se distribuem entre provedor e cliente conforme o serviço.\item provedor é responsável por qualquer configuração do cliente.\item cliente nunca possui responsabilidade de segurança.\item nuvem elimina IAM.\item todos os modelos têm divisão idêntica.\end{enumerate}
\end{questao}

\begin{questao}{S1.80 — Computação em Nuvem}
Uma conta cloud possui storage público por erro de configuração. Esse caso exemplifica:
\begin{enumerate}[label=\Alph*)]\item risco de misconfiguration sob responsabilidade do consumidor conforme o modelo.\item falha física obrigatória do provedor.\item impossibilidade de controle de acesso em nuvem.\item CAP.\item RAID 0.\end{enumerate}
\end{questao}

\begin{questao}{S1.81 — Computação em Nuvem}
FinOps busca:
\begin{enumerate}[label=\Alph*)]\item integrar engenharia, finanças e negócio para otimizar valor/custo da nuvem.\item eliminar monitoramento de custo.\item usar apenas CAPEX.\item substituir segurança.\item impedir autoscaling.\end{enumerate}
\end{questao}

\begin{questao}{S1.82 — Computação em Nuvem}
Estratégia multi-region pode melhorar resiliência, mas:
\begin{enumerate}[label=\Alph*)]\item exige análise de consistência, custo, replicação e failover.\item elimina todos os riscos.\item torna backup desnecessário.\item dispensa teste de recuperação.\item impede latência.\end{enumerate}
\end{questao}

\begin{questao}{S1.83 — Análise de Dados}
A mediana é preferível à média, em certos casos, porque:
\begin{enumerate}[label=\Alph*)]\item é menos sensível a valores extremos.\item sempre usa todos os valores de forma linear.\item não depende de ordenação.\item é igual à média em qualquer distribuição.\item mede dispersão.\end{enumerate}
\end{questao}

\begin{questao}{S1.84 — Análise de Dados}
Em classificação binária muito desbalanceada, acurácia pode ser enganosa porque:
\begin{enumerate}[label=\Alph*)]\item prever majoritariamente a classe dominante pode produzir alta acurácia e baixa detecção da classe rara.\item acurácia é sempre zero.\item recall e precisão são idênticos.\item matriz de confusão não se aplica.\item ROC é impossível.\end{enumerate}
\end{questao}

\begin{questao}{S1.85 — Análise de Dados}
Precision é calculada como:
\begin{enumerate}[label=\Alph*)]\item $TP/(TP+FP)$.\item $TP/(TP+FN)$.\item $TN/(TN+FP)$.\item $(TP+TN)/N$ apenas.\item $FP/(TP+FP)$.\end{enumerate}
\end{questao}

\begin{questao}{S1.86 — Análise de Dados}
Data leakage ocorre quando:
\begin{enumerate}[label=\Alph*)]\item informação indisponível no momento real de previsão contamina treino/validação.\item dados são criptografados.\item há divisão treino/teste correta.\item feature é padronizada dentro do pipeline.\item modelo usa regularização.\end{enumerate}
\end{questao}

\begin{questao}{S1.87 — Análise de Dados}
Correlação alta entre duas variáveis:
\begin{enumerate}[label=\Alph*)]\item não prova causalidade por si só.\item prova necessariamente causalidade.\item significa igualdade das variáveis.\item elimina confundimento.\item substitui experimento.\end{enumerate}
\end{questao}

\begin{questao}{S1.88 — Análise de Dados}
Para avaliar modelo em dados temporais, é geralmente inadequado:
\begin{enumerate}[label=\Alph*)]\item embaralhar aleatoriamente observações futuras com passadas sem respeitar ordem temporal.\item usar janela temporal de validação.\item simular previsão futura.\item evitar leakage temporal.\item monitorar drift.\end{enumerate}
\end{questao}

\begin{questao}{S1.89 — Inteligência Artificial}
RAG é arquitetura em que:
\begin{enumerate}[label=\Alph*)]\item documentos relevantes são recuperados e fornecidos como contexto à geração.\item modelo é treinado do zero em cada pergunta.\item não existe retrieval.\item embeddings são proibidos.\item somente regras fixas são utilizadas.\end{enumerate}
\end{questao}

\begin{questao}{S1.90 — Inteligência Artificial}
Fine-tuning é mais adequado que RAG quando o objetivo principal é:
\begin{enumerate}[label=\Alph*)]\item adaptar comportamento/estilo/padrões do modelo, não apenas consultar base factual mutável.\item atualizar diariamente leis sem retreino.\item citar documentos externos em tempo real.\item aplicar ACL por documento no retrieval.\item substituir qualquer governança de dados.\end{enumerate}
\end{questao}

\begin{questao}{S1.91 — Inteligência Artificial}
Prompt injection em sistema RAG pode ocorrer quando:
\begin{enumerate}[label=\Alph*)]\item conteúdo recuperado contém instruções maliciosas capazes de influenciar o modelo.\item banco usa índice vetorial.\item usuário possui MFA.\item modelo usa tokenização.\item existe cache.\end{enumerate}
\end{questao}

\begin{questao}{S1.92 — Inteligência Artificial}
Concept drift representa:
\begin{enumerate}[label=\Alph*)]\item mudança na relação entre entradas e alvo ao longo do tempo.\item aumento de memória RAM.\item compressão de embedding.\item atualização do sistema operacional.\item normalização de dados estática.\end{enumerate}
\end{questao}

\begin{questao}{S1.93 — Inteligência Artificial}
Um agente de IA com capacidade de executar pagamentos deve possuir, entre outros controles:
\begin{enumerate}[label=\Alph*)]\item limites de ação, autorização, segregação, logs e aprovação humana para ações críticas.\item acesso root irrestrito.\item ausência de registro.\item credencial compartilhada.\item autoaprovação de qualquer valor.\end{enumerate}
\end{questao}

\begin{questao}{S1.94 — Inteligência Artificial}
Explicabilidade em sistema de apoio à decisão é útil para:
\begin{enumerate}[label=\Alph*)]\item compreender fatores da decisão, validar comportamento e apoiar accountability, conforme técnica aplicável.\item garantir que o modelo nunca erre.\item substituir avaliação de viés.\item eliminar necessidade de dados de qualidade.\item provar causalidade automaticamente.\end{enumerate}
\end{questao}

\begin{questao}{S1.95 — Auditoria do Setor Público}
Evidência suficiente e apropriada significa, respectivamente:
\begin{enumerate}[label=\Alph*)]\item quantidade adequada e qualidade/relevância/confiabilidade.\item opinião e volume.\item apenas documentos impressos.\item materialidade e risco.\item causa e efeito.\end{enumerate}
\end{questao}

\begin{questao}{S1.96 — Auditoria do Setor Público}
Materialidade em auditoria:
\begin{enumerate}[label=\Alph*)]\item pode envolver dimensão quantitativa e qualitativa.\item é somente valor financeiro fixo.\item não afeta planejamento.\item elimina risco.\item é igual a amostra.\end{enumerate}
\end{questao}

\begin{questao}{S1.97 — Auditoria do Setor Público}
Um achado de auditoria bem estruturado relaciona:
\begin{enumerate}[label=\Alph*)]\item critério, condição, causa, efeito/risco e evidência.\item somente opinião do auditor.\item apenas recomendação.\item apenas norma.\item apenas valor financeiro.\end{enumerate}
\end{questao}

\begin{questao}{S1.98 — Auditoria do Setor Público}
Amostragem em auditoria:
\begin{enumerate}[label=\Alph*)]\item deve ser desenhada de acordo com objetivo, população, risco e inferência pretendida.\item pode sempre ser escolhida por conveniência sem documentação.\item elimina risco de amostragem.\item exige examinar 100\% dos itens.\item é incompatível com análise de dados.\end{enumerate}
\end{questao}

\begin{questao}{S1.99 — Auditoria do Setor Público}
Independência e objetividade são relevantes porque:
\begin{enumerate}[label=\Alph*)]\item reduzem risco de julgamento influenciado por interesses indevidos.\item impedem qualquer comunicação com auditado.\item eliminam necessidade de evidência.\item substituem competência técnica.\item autorizam preconceito profissional.\end{enumerate}
\end{questao}

\begin{questao}{S1.100 — Auditoria do Setor Público}
Monitoramento de recomendação busca verificar:
\begin{enumerate}[label=\Alph*)]\item implementação e efeitos das providências acordadas/determinadas, conforme o caso.\item apenas se o relatório foi publicado.\item se a equipe mudou.\item quantidade de páginas.\item se houve nova contratação de auditor.\end{enumerate}
\end{questao}


\clearpage
\section{Gabarito comentado — Simulado 1}

\begin{longtable}{c c p{10.8cm}}
\toprule
Q & Gab. & Comentário \\
\midrule
\endfirsthead
\toprule Q & Gab. & Comentário \\
\midrule
\endhead
1 & B & ``Embora'' expressa concessão; ``apesar de'' preserva a relação e o resultado negativo. \\
2 & B & Oração relativa restritiva, sem vírgulas, delimita o subconjunto devolvido. \\
3 & C & ``Existir'' concorda com ``evidências''; haver/fazer impessoais ficam no singular. \\
4 & B & ``Porquanto'' pode assumir valor causal/explicativo próximo de ``porque''. \\
5 & C & ``Seu'' pode retomar equipe ou gestor sem contexto desambiguador. \\
6 & B & ``Referência a'' + artigo feminino ``a norma'' produz crase. \\
7 & C & Validar domínio por canal confiável e reportar phishing evita captura de credenciais. \\
8 & B & Necessidade = tratamento limitado ao mínimo necessário à finalidade. \\
9 & B & A LAI adota publicidade como regra e sigilo como exceção legal. \\
10 & B & Privacidade, segurança e finalidade continuam exigíveis em serviços digitais. \\
11 & B & Negação de universal é existencial: existe ao menos um que não foi aprovado. \\
12 & A & Contrapositiva de $P\rightarrow Q$ é $\neg Q\rightarrow\neg P$. \\
13 & C & $800\times1{,}15=920$ mil. \\
14 & C & Inclusão-exclusão: $120+80-40=160$. \\
15 & C & De $x-y=2$, $x=y+2$; substituindo: $3y+2=14$, $y=4$, $x=6$. \\
16 & B & Desconcentração distribui competência dentro da mesma pessoa jurídica. \\
17 & C & Autoexecutoriedade permite execução direta em hipóteses admitidas pelo ordenamento. \\
18 & B & Anulação decorre de ilegalidade; revogação, em regra, de mérito administrativo. \\
19 & B & A disciplina vigente da improbidade exige dolo e enquadramento legal; ilegalidade isolada não basta. \\
20 & B & Regra constitucional: responsabilidade objetiva do Estado perante a vítima, com análise de nexo/excludentes. \\
21 & B & Municípios integram a organização federativa brasileira e possuem autonomia constitucional. \\
22 & B & Direitos fundamentais não se restringem ao art. 5º. \\
23 & B & O Congresso exerce controle externo com auxílio do TCU. \\
24 & B & CPI não pratica atos submetidos à reserva jurisdicional. \\
25 & B & MP é função essencial à Justiça. \\
26 & B & Controle interno e externo têm posições e competências distintas, mas complementares. \\
27 & B & Auditoria operacional examina desempenho, economicidade, eficiência, eficácia e efetividade. \\
28 & B & Jurisdição una assegura acesso judicial contra lesão ou ameaça a direito. \\
29 & B & Parecer prévio subsidia o julgamento político competente das contas de governo. \\
30 & A & A IN nº 50/2017 disciplina tomada de contas especial no TCE-MA. \\
31 & A & A IN nº 82/2025 trata de emendas parlamentares, transparência e rastreabilidade. \\
32 & B & Regimento não pode contrariar a Lei Orgânica. \\
33 & B & Rastreabilidade é reconstrução da cadeia completa da emenda e sua execução. \\
34 & A & Guaxenduba pertence ao conflito contra os franceses. \\
35 & A & Balaiada: Período Regencial, forte participação popular e causas sociais/políticas. \\
36 & B & Parnaíba, Gurupi e Tocantins-Araguaia são destacados como bacias limítrofes. \\
37 & A & Itaqui e corredores ferroviários integram cadeias nacionais e internacionais de exportação. \\
38 & B & A LBI adota modelo relacional: impedimento em interação com barreiras. \\
39 & A & ODS 16 trata de paz, justiça e instituições eficazes, responsáveis e inclusivas. \\
40 & A & Ações afirmativas buscam corrigir desigualdades e promover igualdade material. \\
41 & D & Em /26 restam 6 bits: $2^6=64$ endereços no bloco. \\
42 & A & 802.1X + EAP + RADIUS é arquitetura típica de controle de acesso à rede. \\
43 & A & Estado preso ao nó quebra failover; externalização/replicação ou desenho stateless reduz o problema. \\
44 & A & RAID e snapshot local compartilham riscos e não cobrem exclusão, ransomware e desastre do storage. \\
45 & A & São as quatro condições clássicas de Coffman para deadlock. \\
46 & B & AD usa LDAP, Kerberos, DNS e outros componentes; não é sinônimo de LDAP. \\
47 & A & Dependência transitiva de atributo não chave viola 3FN. \\
48 & A & MVCC mantém versões para coordenar concorrência e reduzir bloqueios conforme implementação. \\
49 & A & Fatos registram medidas/eventos e referenciam dimensões. \\
50 & A & CDC captura mudanças para integração/streaming. \\
51 & A & Lineage documenta origem, transformações e destino. \\
52 & A & Sob partição, CAP explicita trade-off C versus A. \\
53 & A & Requisito não mensurável dificulta teste e aceite. \\
54 & A & Sprint Review inspeciona incremento e adapta backlog com stakeholders. \\
55 & A & Ciclo Red-Green-Refactor. \\
56 & A & Credencial admin permanente viola menor privilégio e segregação. \\
57 & A & Readiness indica aptidão para receber tráfego. \\
58 & A & SLO precisa de SLI/métrica e janela de avaliação. \\
59 & A & MFA fatigue explora repetição de prompts para induzir aprovação. \\
60 & A & PAM controla contas/acessos privilegiados. \\
61 & A & Assinatura digital apoia integridade, autenticidade e não repúdio; não cifra automaticamente conteúdo. \\
62 & A & Cópia imutável/offline reduz comprometimento simultâneo por ransomware. \\
63 & A & Cadeia de custódia documenta percurso e integridade da evidência. \\
64 & A & Zero Trust = verificar continuamente, contexto, identidade e menor privilégio. \\
65 & A & EDM é domínio de governança no COBIT. \\
66 & A & Governança avalia, direciona e monitora; gestão planeja/constrói/executa/monitora atividades. \\
67 & A & ITIL 4 enfatiza cocriar valor em relações de serviço. \\
68 & A & Probabilidade e impacto são componentes centrais da análise de risco. \\
69 & A & Gateways controlam ramificação e convergência de fluxos BPMN. \\
70 & A & KPI local pode incentivar comportamento desalinhado do resultado final. \\
71 & A & Essas são as três fases macro da IN SGD/ME 94/2022. \\
72 & A & Ignorar referências públicas comparáveis pode distorcer preço estimado e elevar risco de sobrepreço. \\
73 & A & Evidência produzida apenas pela parte interessada é pouco independente. \\
74 & A & Portabilidade, interoperabilidade e saída testada mitigam lock-in. \\
75 & B & 30 dias = 43.200 min; 0,1\% = 43,2 min. \\
76 & A & Dimensionar pelo total sem relação com demanda gera licenças ociosas e sobrecontratação. \\
77 & A & Em IaaS, cliente normalmente gerencia SO, middleware/aplicações e configurações sob sua responsabilidade. \\
78 & A & Elasticidade é ajuste dinâmico de capacidade. \\
79 & A & A divisão varia com IaaS/PaaS/SaaS e configuração concreta. \\
80 & A & Storage público por configuração é exemplo clássico de responsabilidade do consumidor. \\
81 & A & FinOps conecta custo, valor e decisões técnicas/financeiras. \\
82 & A & Multi-region aumenta complexidade de dados, consistência, custo e failover; exige teste. \\
83 & A & Mediana sofre menos influência de outliers. \\
84 & A & Classe dominante pode inflar acurácia; examine precision/recall e matriz de confusão. \\
85 & A & Precision = verdadeiros positivos entre todos os positivos previstos. \\
86 & A & Leakage usa informação que não estaria disponível no momento real da previsão. \\
87 & A & Correlação não demonstra causalidade. \\
88 & A & Em série temporal, mistura aleatória futuro-passado pode vazar informação. \\
89 & A & RAG recupera contexto externo e o fornece ao modelo gerador. \\
90 & A & Fine-tuning é apropriado para comportamento/padrões; RAG é melhor para conhecimento factual atualizável e rastreável. \\
91 & A & Conteúdo recuperado pode carregar instruções maliciosas e contaminar a geração. \\
92 & A & Concept drift altera a relação entre características e alvo. \\
93 & A & Ações com impacto financeiro exigem limites, autorização, logs e aprovação/segregação adequados. \\
94 & A & Explicabilidade apoia validação, entendimento e accountability, sem garantir perfeição. \\
95 & A & Suficiência = quantidade; adequação = relevância e confiabilidade/qualidade. \\
96 & A & Materialidade pode ser financeira e qualitativa, afetando escopo e julgamento. \\
97 & A & Essa cadeia estrutura um achado auditável. \\
98 & A & Amostra depende do objetivo e da população; conveniência não é regra geral. \\
99 & A & Independência e objetividade protegem julgamento profissional contra influência indevida. \\
100 & A & Monitoramento verifica implementação e efeito das providências após a auditoria. \\
\bottomrule
\end{longtable}

\subsection*{Folha de resultado}

\begin{itemize}
\item Gerais: \underline{\hspace{2cm}}/40.
\item Específicas: \underline{\hspace{2cm}}/60.
\item Nota ponderada: $G+2E=$ \underline{\hspace{2.5cm}}/160.
\item Disciplinas com desempenho abaixo de 70\%: \underline{\hspace{8cm}}.
\item Erros por interpretação: \underline{\hspace{1.5cm}}; teoria: \underline{\hspace{1.5cm}}; cálculo: \underline{\hspace{1.5cm}}; distração: \underline{\hspace{1.5cm}}.
\end{itemize}


\section{Métrica de desempenho}

Calcule:
\[
N = G + 2E,
\]
onde $G$ é o número de acertos em conhecimentos gerais e $E$ o número de acertos em conhecimentos específicos.

Além da nota, registre desempenho por eixo. Um candidato que acerta 80\% no total mas erra repetidamente Segurança, Contratos ou Auditoria ainda possui risco alto, porque questões específicas valem o dobro e alimentam a discursiva.

\begin{resumo}[meta de treinamento]
Use três metas progressivas: \textbf{70\%} para consolidação; \textbf{80\%} para competitividade; \textbf{85\%+} com baixa variância entre disciplinas para reta final. Essas metas são de estudo, não notas de corte oficiais.
\end{resumo}
